    \documentclass[14pt,a4paper]{extarticle}

    \usepackage[utf8]{inputenc}
    \usepackage[english]{babel}
    \usepackage{setspace}
    \usepackage{geometry}

    \geometry{
        left=25mm,
        right=15mm,
        top=20mm,
        bottom=20mm
    }

    \setstretch{1.5}

    \begin{document}

    This course paper investigates the acceleration of genetic algorithms through CUDA-based GPU parallelization, covering both performance analysis and the development of reusable implementation components.

    \bigskip\textbf{Hypothesis.} The application of CUDA technology to parallelize computations in genetic algorithms 
    will enable significant acceleration of key operations: fitness function evaluation, 
    selection, crossover, and mutation - compared to sequential implementation on the CPU. 
    It is expected that parallelization efficiency will increase with population
    size and computational complexity of the fitness function, subject to the theoretical limits imposed by Amdahl's Law and data transfer overhead.

    \bigskip Parallel processing of large populations on the GPU is expected to improve the quality of
    solutions found, because it allows operation with significantly larger population sizes without
    a proportional increase in overall algorithm execution time. This opens up new
    possibilities for solving complex optimization problems in acceptable timeframes.

    \bigskip\textbf{Object of the Research.} Genetic algorithms represent the general domain of the research paper, serving as black-box optimization material.

    \bigskip\textbf{Subject of the Research.} Parallelization strategies and speedups achieved through CUDA form the specific focus, targeting custom kernel development for selection, crossover, mutation, and fitness evaluation on GPU.
    ​

    \bigskip\textbf{Research Gap.} A comparative analysis of publication volumes on Semantic Scholar\footnote{Semantic Scholar Academic Graph API v1. Endpoint: \texttt{GET .../graph/v1/paper/\allowbreak search?\allowbreak query=\{Q\}\&limit=1} (base URL: \texttt{api.semanticscholar.org}). The \texttt{total} field in the JSON response gives the match count; \texttt{\&year=2020-2025} was appended for five-year filtering. All queries: October~2025. Verification in March~2026: control returned \texttt{total:\,286918} (all-time) and \texttt{total:\,109069} (5-year), confirming proportionality with the original data.} (accessed October 2025) shows that GPU-accelerated genetic algorithm optimization is substantially less studied than established optimization domains. Table~\ref{tab:ss-gap} summarizes the result counts.

    \begin{table}[h]
    \centering
    \caption{Semantic Scholar publication counts by query (October 2025).}
    \label{tab:ss-gap}
    \small
    \begin{tabular}{p{0.55\textwidth}rrr}
    \hline
    \textbf{Query} & \textbf{All-time} & \textbf{5-year} & \textbf{Ratio}\textsuperscript{*} \\
    \hline
    databases clustering optimization (control) & 243{,}000 & 94{,}500 & 1$\times$ \\
    CUDA parallel processing genetic algorithms & 8{,}800 & 1{,}380 & 28$\times$ \\
    MapReduce parallel processing genetic algorithms optimization & 7{,}530 & 1{,}150 & 32$\times$ \\
    OpenCL parallel processing genetic algorithms & 6{,}250 & 946 & 39$\times$ \\
    \hline
    \multicolumn{4}{l}{\textsuperscript{*}Ratio = control all-time / query all-time, rounded.}
    \end{tabular}
    \end{table}

    Within this space, the majority of GPU implementations use native C/C++ CUDA kernels or high-level framework abstractions (EvoJAX, EvoTorch, EvoX). Python-native CUDA kernel development via Numba, despite NVIDIA's investment in a Python CUDA SDK, appears rarely in the evolutionary computation literature. This gap motivates an investigation into whether Python-native GPU kernel development can achieve competitive speedups while retaining the prototyping speed and accessibility advantages that Python offers over lower-level CUDA C/C++ workflows.

    \bigskip\textbf{Structure of the Paper.} Section~2 provides the theoretical background: a literature review of GPU-accelerated genetic algorithms (Section~2.1) and an overview of CUDA architecture with kernel design principles (Section~2.2). Section~3 presents the experimental methodology, benchmark configurations, and results for two test problems: XOR perceptron training (Section~3.2) and the Traveling Salesman Problem (Section~3.3), followed by a cross-benchmark analysis (Section~3.4). Section~4 describes the reusable GPU-accelerated GA library developed from the experimental code and validates its results against the original benchmarks. Section~5 summarizes the findings and outlines directions for future work.

    \end{document}
